\documentclass[10pt,letterpaper]{article}
\usepackage{cornell-notes}

\title{Chapter 9: Reading List and Bibliography}
\author{}
\date{}

\setDocTitle{Chapter 9: Reading List and Bibliography}
\setDocSubtitle{Computer Networks}
\setDocVersion{v1.0}
\setDocStatus{Study Companion}
\setDocOwner{Computer Networks Cornell Notes}
\setDocAuthor{}
\setDocDate{\today}
\setDocSummary{Cornell-format study and review notes for Chapter 9: Reading List and Bibliography.}

\setCornellCollection{Computer Networks Cornell Notes}
\setCornellUnitType{Chapter}
\setCornellUnitNumber{9}
\setCornellUnitTitle{Reading List and Bibliography}
\setCornellCompanionLabel{Study and review companion}

\hypersetup{pdftitle={Chapter 9: Reading List and Bibliography},pdfsubject={Cornell notes},pdfauthor={}}

\begin{document}
\maketitle
\notesection{Chapter overview}
\begin{overviewbox}
\textbf{Learning objectives}\begin{itemize}
  \item Select further reading by technical layer, learning goal, and evidence type.
  \item Distinguish tutorials, surveys, standards, histories, measurement studies, implementation texts, and research papers.
  \item Map the recommended reading themes to Chapters 1 through 8.
  \item Use the alphabetical bibliography as a lookup index without treating every citation as an annotated recommendation.
\end{itemize}
\textbf{Prerequisite concepts}\quad The concepts, protocols, and layer relationships developed in Chapters 1 through 8.\par
\textbf{Key themes}\quad Deliberate study sequencing; primary versus secondary evidence; history; measurement; implementation; research venues; bibliographic lookup.\par
\textbf{Named mechanisms and standards}\quad Layer-keyed reading paths, journals and magazines, conference proceedings, standards and RFCs, surveys, textbooks, implementation studies, and an alphabetical reference index.
\end{overviewbox}

\notesection{Cornell cue-and-note record}

\begin{longtable}{@{}L{0.245\textwidth}|L{0.695\textwidth}@{}}
\rowcolor{CNavy}\color{white}\bfseries Cue / recall prompt &
\color{white}\bfseries Notes \\
\endfirsthead
\rowcolor{CNavy}\color{white}\bfseries Cue / recall prompt &
\color{white}\bfseries Notes (continued) \\
\endhead
\hline
\endfoot
\cellcolor{CCue}\textbf{9.1 Suggestions for further reading} & \begin{minipage}[t]{\linewidth}\raggedright The suggested pathway begins with peer-reviewed networking journals, society magazines, specialist group publications, and conferences. Journals emphasize sustained research results; magazines often provide surveys and case studies; conferences expose emerging systems and measurements. Full reference details are resolved through Section 9.2.\end{minipage} \\ \hline
\cellcolor{CCue}\textbf{9.1.1 Introduction and general works} & \begin{minipage}[t]{\linewidth}\raggedright Representative directions include Comer's introductory work, historical collections, Crovella and Krishnamurthy on measurement, and histories of early Internet development.\par\begin{itemize}\item \textbf{Accessible orientation} Use broad introductions to consolidate Internet history, architecture, protocols, growth, and services before specializing.\item \textbf{Historical primary material} Anniversary collections and histories reveal how TCP, multicast, DNS, Ethernet, packet switching, and architectural decisions developed.\item \textbf{Measurement} Internet measurement literature asks how to infer topology, performance, and application behavior in a decentralized system.\item \textbf{Standards process} Critical accounts of standards adoption show that interoperability work is also shaped by patents, vendors, and governance.\end{itemize}\end{minipage} \\ \hline
\cellcolor{CCue}\textbf{9.1.2 Physical layer} & \begin{minipage}[t]{\linewidth}\raggedright The list moves from general digital telephony to media-specific depth and historical context.\par\begin{itemize}\item \textbf{Telephony and multiplexing} Study digital transmission, switching, fiber, cellular access, and DSL as an integrated carrier system.\item \textbf{Satellite Internet} Extend propagation-delay concepts into satellite-specific routing and switching concerns.\item \textbf{Telecommunications history} Use historical surveys to connect telegraphy, radio, switching, cable, optical systems, mobile phones, packet switching, and the Internet.\item \textbf{Fiber, cellular, and RFID depth} Specialist texts expand optical devices and waveguides, 3G radio interfaces, and passive RFID behavior beyond the chapter overview.\end{itemize}\end{minipage} \\ \hline
\cellcolor{CCue}\textbf{9.1.3 Data link layer} & \begin{minipage}[t]{\linewidth}\raggedright The reading path separates mathematical coding depth from link-protocol engineering.\par\begin{itemize}\item \textbf{Carrier Ethernet} Examine how Ethernet framing and operations are extended beyond a local-area setting.\item \textbf{Error-control coding} Develop the algebra and decoding theory behind Hamming, low-density parity-check, and related codes.\item \textbf{Link reliability} Use broader data-communications treatments to revisit detection, retransmission, and flow control as a connected design space.\end{itemize}\end{minipage} \\ \hline
\cellcolor{CCue}\textbf{9.1.4 MAC sublayer} & \begin{minipage}[t]{\linewidth}\raggedright \begin{itemize}\item \textbf{Broadband wireless} A WiMAX-focused treatment deepens OFDM, multiple antennas, and scheduled multi-access.\item \textbf{Wireless LANs} An IEEE 802.11 guide expands MAC, physical-layer variants, and security, while its publication date limits coverage of later amendments.\item \textbf{Bridges and routing} Perlman's treatment connects bridge design, spanning trees, routers, and general protocol reasoning from an implementer's perspective.\end{itemize}\end{minipage} \\ \hline
\cellcolor{CCue}\textbf{9.1.5 Network layer} & \begin{minipage}[t]{\linewidth}\raggedright \begin{itemize}\item \textbf{Internet protocol suite} Protocol-suite references deepen IP and adjacent control protocols through detailed examples.\item \textbf{Mobile IP networks} Specialized works link IP design, IPv6, mobility, and operator-network engineering.\item \textbf{Routing protocols} Routing texts deepen RIP, CIDR, OSPF, BGP, unicast, and multicast, while older publication dates require careful historical labeling.\item \textbf{Operations and measurement} Large-network design works add traffic engineering, service delivery, deployment, and management.\item \textbf{Router implementation} Network algorithmics shifts attention from distributed protocol exchange to high-speed lookup, classification, scheduling, and forwarding inside a router.\end{itemize}\end{minipage} \\ \hline
\cellcolor{CCue}\textbf{9.1.6 Transport layer} & \begin{minipage}[t]{\linewidth}\raggedright \begin{itemize}\item \textbf{TCP and UDP depth} Detailed protocol references and illustrated traces deepen header, state, timer, and congestion behavior.\item \textbf{Challenged networks} Delay- and disruption-tolerant networking literature expands store-carry-forward architecture, protocols, and applications where no continuous path exists.\end{itemize}\end{minipage} \\ \hline
\cellcolor{CCue}\textbf{9.1.7 Application layer} & \begin{minipage}[t]{\linewidth}\raggedright \begin{itemize}\item \textbf{Web origins and architecture} Historical Web material explains URLs, HTTP, HTML, and early architectural intent.\item \textbf{Content delivery} CDN literature focuses on real placement, operation, caching, and performance decisions.\item \textbf{Structured Web data and services} XML-oriented references expand representation and then-current service techniques.\item \textbf{Web protocols and measurement} Comprehensive Web treatments connect clients, servers, proxies, caching, traffic, and research.\item \textbf{Audio/video applications} Video-over-IP and voice-over-IP works bridge media engineering, quality of service, and signaling protocols.\end{itemize}\end{minipage} \\ \hline
\cellcolor{CCue}\textbf{9.1.8 Network security} & \begin{minipage}[t]{\linewidth}\raggedright \begin{itemize}\item \textbf{Security engineering} System-oriented reading connects cryptography with people, operations, physical context, and real application failures.\item \textbf{Cryptographic protocol design} Engineering-focused work explains why primitives are composed in particular ways and how misuse defeats goals.\item \textbf{Steganography} Specialist literature studies hiding and detecting information within digital media.\item \textbf{Protocol depth} Technical references expand secret/public-key algorithms, hashes, authentication, Kerberos, PKI, IPsec, SSL/TLS, and mail protection.\item \textbf{Adversarial and human views} Broader and attacker-oriented works show why secure design must account for environments, users, and system-wide attack paths.\end{itemize}\end{minipage} \\ \hline
\cellcolor{CCue}\textbf{9.2 Alphabetical bibliography} & \begin{minipage}[t]{\linewidth}\raggedright This section is a reference-resolution index, not a set of uniform endorsements or annotations. Use author/title keys from Section 9.1 to locate complete publication details. Unannotated entries may support specific claims elsewhere, but no relevance statement should be invented from citation metadata alone.\end{minipage} \\ \hline
\end{longtable}

\begin{legacybox}Many recommended works describe technologies and predictions current at their publication dates, including 3G, WiMAX, older 802.11 amendments, early Web-service techniques, and older routing or security mechanisms. Read them for foundations and history, and verify present-day deployment or security guidance separately.\end{legacybox}

\notesection{Terminology and notation}

\begin{longtable}{@{}L{0.245\textwidth}|L{0.695\textwidth}@{}}
\rowcolor{CNavy}\color{white}\bfseries Cue / recall prompt &
\color{white}\bfseries Notes \\
\endfirsthead
\rowcolor{CNavy}\color{white}\bfseries Cue / recall prompt &
\color{white}\bfseries Notes (continued) \\
\endhead
\hline
\endfoot
\cellcolor{CCue}\textbf{Primary source} & \begin{minipage}[t]{\linewidth}\raggedright Original specification, experiment, measurement, design paper, or firsthand historical account.\end{minipage} \\ \hline
\cellcolor{CCue}\textbf{Secondary source} & \begin{minipage}[t]{\linewidth}\raggedright Tutorial, survey, textbook, or synthesis that interprets multiple primary sources.\end{minipage} \\ \hline
\cellcolor{CCue}\textbf{Survey} & \begin{minipage}[t]{\linewidth}\raggedright Structured overview that organizes a research area and its major approaches.\end{minipage} \\ \hline
\cellcolor{CCue}\textbf{Case study} & \begin{minipage}[t]{\linewidth}\raggedright Detailed analysis of one implementation, deployment, incident, or operational setting.\end{minipage} \\ \hline
\cellcolor{CCue}\textbf{Proceedings} & \begin{minipage}[t]{\linewidth}\raggedright Published collection of papers presented at a conference or symposium.\end{minipage} \\ \hline
\cellcolor{CCue}\textbf{RFC} & \begin{minipage}[t]{\linewidth}\raggedright Numbered Internet document whose status may be standards-track, informational, experimental, historic, or otherwise classified.\end{minipage} \\ \hline
\cellcolor{CCue}\textbf{Annotated recommendation} & \begin{minipage}[t]{\linewidth}\raggedright Citation accompanied by a stated reason, scope, audience, or limitation.\end{minipage} \\ \hline
\cellcolor{CCue}\textbf{Bibliographic index} & \begin{minipage}[t]{\linewidth}\raggedright Lookup structure containing publication metadata without necessarily explaining relevance.\end{minipage} \\ \hline
\end{longtable}

\notesection{Comparison matrix}

\begin{longtable}{@{}L{0.313\textwidth}L{0.313\textwidth}L{0.313\textwidth}@{}}
\toprule
\textbf{Evidence type} & \textbf{Best use} & \textbf{Reading caution} \\ \midrule
Textbook / monograph & Build a coherent conceptual model & May lag fast-changing standards or deployment \\
Survey / tutorial & Orient quickly and discover terminology & Selection and synthesis reflect author scope \\
Standard / RFC & Resolve normative protocol behavior & Status, updates, and errata must be checked \\
Research paper & Study a precise contribution and evidence & Assumptions may be narrow; later work may revise conclusions \\
Measurement study & Understand observed behavior & Results depend on vantage point, time, workload, and method \\
Historical account & Understand design context and evolution & Not a current implementation or security guide \\
\bottomrule
\end{longtable}
\notesection{Chapter synthesis}

\begin{overviewbox}\textbf{Study path.} Begin with an accessible synthesis, then read a survey to map the area, a standard to resolve protocol requirements, an implementation or measurement study to see operational behavior, and selected research papers to examine unresolved tradeoffs.\par
\textbf{Common confusions.} A bibliography entry is not automatically a recommendation; an RFC is not automatically a current Internet Standard; an older authoritative text can remain conceptually valuable while its deployment details age.\par
\textbf{Derived insight.} The layer-keyed reading list is most useful as a graph rather than a queue: follow prerequisite links and a concrete question instead of reading every citation alphabetically.\end{overviewbox}

\notesection{Retrieval practice}

\begin{enumerate}
  \item Which publication type would you choose first to learn an unfamiliar networking subfield, and why?
  \item How would you validate a protocol detail found in an older textbook?
  \item Design a three-source reading path for link-state routing using a survey, a specification, and an implementation study.
  \item Why are Internet measurement conclusions sensitive to vantage point and time?
  \item Which readings connect Chapters 2 and 4 through wireless physical and multi-access design?
  \item How does network algorithmics extend the Chapter 5 view of routing and forwarding?
  \item Why should older security or wireless recommendations be labeled as historical context?
  \item How should Section 9.2 be used when a citation has no accompanying annotation?
\end{enumerate}
\begin{CornellSummaryBox}{Rapid-review Cornell summary}Use the reading list by question and layer: broad works establish vocabulary; surveys organize a field; standards resolve protocol requirements; implementation and measurement studies expose operational reality; research papers deepen specific tradeoffs. Section 9.2 supplies metadata, not invented relevance. Publication date and document status are part of technical interpretation.\end{CornellSummaryBox}
\end{document}
